% Context from chapters-tex/sec-optim-smooth.tex; for mathematical reference, not compilation.
s a global minimizer.
If $f$ is differentiable and convex, 
\eq{
	x^\star \in \uargmin{x} f(x) 
	\quad\Longleftrightarrow\quad
	\nabla f(x^\star)=0.
}
\end{prop}
\begin{proof}
Fix $x$ and choose $0<t<1$ so that $tx+(1-t)x^\star$ lies in the neighborhood where $x^\star$ is minimal. Local optimality and convexity then give
\eq{
	f(x^\star) \leq f(tx + (1-t)x^\star) \leq t f(x) + (1-t) f(x^\star)
	\qarrq
	f(x^\star) \leq f(x)
}
proving that $x^\star$ is a global minimizer.

For the second part, we already saw in Proposition~\ref{prop-cs-min} the $\Rightarrow$ implication. We assume that $\nabla f(x^\star)=0$. Since the graph of $f$ is above its tangent by convexity (as stated in Proposition~\ref{prop-above-tgt}),
\eq{
	f(x) \geq f(x^\star) + \dotp{\nabla f(x^\star)}{x-x^\star} = f(x^\star).
}
\end{proof}

For a differentiable convex objective, minimization is therefore equivalent to solving $\nabla f(x)=0$, a system of $p$ equations in $p$ unknowns.
 
An explicit solution is often unavailable, but the equation still characterizes the minimizers $x^\star$.



                                                                                  
\subsection{Least Squares}

The gradient of the least-squares objective~\eqref{eq-least-square} follows by expanding the squared norm:
\begin{align*}
	f(x+\epsilon) &= \frac{1}{2}\norm{Ax-y+A\epsilon}^2 = \frac12\norm{Ax-y}^2+\dotp{Ax-y}{A\epsilon} + \frac{1}{2}\norm{A\epsilon}^2 \\
		&=f(x) + \dotp{\epsilon}{A^\top(Ax-y)} + o(\norm{\epsilon}).
\end{align*} 
We used $\norm{A\epsilon}^2 = o(\norm{\epsilon})$ and the transpose $A^\top$.
 
The transpose exchanges rows and columns, $A^\top = (A_{j,i})_{i=1,\ldots,n}^{j=1,\ldots,p}$. Its essential role in gradient calculations is the adjoint identity
\eq{
	\foralls (u,v) \in \RR^{p} \times \RR^n, \quad
	\dotp{A u}{v}_{\RR^n} = \dotp{u}{A^\top v}_{\RR^p}. 
}
Computing gradients of functions involving linear operators requires such a transposition step.
 
This computation shows that
\eql{\label{eq-grad-ls}
	\nabla f(x) = A^\top (A x - y). 
}
Hence every minimizer $x^\star$ of $f(x)$ satisfies the normal equations $(A^\top A) x^\star = A^\top y$.
 
If $A^\star A \in \RR^{p \times p}$ is invertible, then $f$ has a single minimizer, namely  
\eql{\label{eq-sol-leastsquare}
	x^\star = (A^\top A)^{-1} A^\top y.
}
Thus $x^\star$ depends linearly on $y$. The operator $(A^\top A)^{-1} A^\star$ is the Moore--Penrose pseudoinverse of $A$; an ordinary inverse is unavailable when $p \neq n$.
 
The condition that $A^\top A$ is invertible is equivalent to $\ker(A)=\{0\}$, since 
\eq{
	A^\top A x = 0 \qarrq \norm{Ax}^2 = \dotp{A^\top A x}{x} = 0 \qarrq A x= 0.
}
When $n<p$, the system is underdetermined and injectivity is impossible. If $n\geq p$ and the entries of $A$ have a joint density, then $\ker(A)=\{0\}$ almost surely. In general, injectivity is equivalent to the feature vectors $a_i$ spanning $\RR^p$.
 
                                                                                  
\subsection{Link with PCA}

Assume the feature vectors are centered, $\sum_i a_i=0$. Otherwise, replace each $a_i$ by $a_i-m$, where $m\eqdef n^{-1}\sum_i a_i$ is the empirical mean. Write $C\eqdef A^\top A$. Then $C/n$ is the empirical covariance matrix of the point cloud. With $a_i=(a_{i,1},\ldots,a_{i,p})^\top$ and $A=(a_{i,j})_{i,j}$, its entries are
 \eq{
 	\foralls (k,\ell) \in \{1,\ldots,p\}^2, \quad
 	\frac{C_{k,\ell}}{n} = \frac{1}{n} \sum_{i=1}^n a_{i,k} a_{i,\ell}. 
 }
 In particular, $C_{k,k}/n$ is the variance along the axis $k$. More generally, for any unit vector $u \in \RR^p$, $\dotp{C u}{u}/n \geq 0$ is the variance along the axis $u$.
 
For instance, in dimension $p=2$, 
\eq{
	\frac{C}{n} = 
	\frac{1}{n}
	\begin{pmatrix}
		 \sum_{i=1}^n a_{i,1}^2 & \sum_{i=1}^n a_{i,1} a_{i,2} \\
		\sum_{i=1}^n a_{i,1}a_{i,2} & \sum_{i=1}^n a_{i,2}^2
	\end{pmatrix}.
}

Since $C$ is symmetric, choose an orthogonal matrix $U=(u_1,\ldots,u_p)$ whose columns are eigenvectors of $C/n$. If $(C/n)u_k=\la_k u_k$, then $C/n=U\diag(\la_k)U^\top$ and $U^{-1}=U^\top$. In these coordinates, the quadratic form is a weighted sum of squares:
\eql{\label{eq-pca-decomp}
	\frac{1}{n}  \dotp{Cx}{x} = \dotp{U \diag(\la_k) U^\top x}{x}
	= \dotp{\diag(\la_k) (U^\top x)}{(U^\top x)}
	= \sum_{k=1}^p \la_k \dotp{x}{u_k}^2.
} 
Here $( U^\top x )_k = \dotp{x}{u_k}$ is the coordinate $k$ of $x$ in the basis $U$. Since $\dotp{Cx}{x}=\norm{Ax}^2$, this shows that all eigenvalues are nonnegative.

\begin{figure}[tbp]
\centering
\BookDrawing{tikz/smooth-pca-quadratic-geometry}
\caption{Left: point clouds $(a_i)_i$ and their principal directions. Right: the quadratic part of $f(x)$.}
\label{fig-link-pca}
\end{figure}

Order the eigenvalues as $\la_1 \geq \la_2 \geq \ldots \geq \la_p$. Projection of $a_i$ onto the leading $m$ eigenvectors gives principal component analysis (PCA), the optimal linear reconstruction in squared error. It is used for compression, dimensionality reduction, and visualization in two dimensions ($m=2$) or three dimensions ($m=3$).

When $C$ is positive definite, its covariance ellipsoid has principal axes $(u_k)_k$ and semiaxes proportional to the standard deviations $\sqrt{\la_k}$.
 
This ellipsoid closely describes a Gaussian cloud with density proportional to $\exp(-\frac n2\dotp{C^{-1}a}{a})$.
 
For the quadratic objective $\frac{1}{2}\dotp{Cx}{x}$ in $f(x)$, the level ellipsoid $\enscond{x}{\frac{1}{n}\dotp{Cx}{x} \leq 1}$ has the same principal axes but reciprocal widths $1/\sqrt{\la_k}$.
 
Figure~\ref{fig-link-pca} shows the reciprocal axis widths in dimension $p=2$. If $C$ is singular, the covariance ellipsoid is degenerate and the quadratic sublevel sets are unbounded along $\ker(C)$.



                                                                                  
\subsection{Classification}

For the classification objective~\eqref{eq-classif}, assume $L$ is differentiable and apply~\eqref{eq-grad-dfn} at $-\diag(y) Ax$:
\begin{align*}
	f(x+\epsilon) &= L( -\diag(y) Ax -\diag(y)  A\epsilon) \\
	&= L(-\diag(y) Ax) + \dotp{\nabla L( -\diag(y) Ax)}{ -\diag(y)  A\epsilon } + o(\norm{\diag(y)  A\epsilon}).
\end{align*} 
Using the fact that $o(\norm{\diag(y)  A\epsilon}) = o(\norm{\epsilon})$, one obtains
\begin{align*}
	f(x+\epsilon) &= f(x) + \dotp{\nabla L( -\diag(y) Ax)}{-\diag(y)  A\epsilon} + o(\norm{\epsilon}) \\
		 &=  f(x) + \dotp{-A^\top \diag(y) \nabla L( -\diag(y) Ax)}{\epsilon} + o(\norm{\epsilon}), 
\end{align*} 	
Using $(AB)^\top = B^\top A^\top$ and $\diag(y)^\top=\diag(y)$ gives
\eq{
	\nabla f(x) = -A^\top \diag(y) \nabla L( -\diag(y) Ax). 
}
Since $L(z)=\sum_i\ell(z_i)$, its gradient is $\nabla L(z)=(\ell'(z_i))_{i=1}^n$. For th